% 5. A Synthesis: Two Modes, One System
\section{A Synthesis: Two Modes, One System}
\label{sec:synthesis}

If I put the case for and the case against side by side, I do not think they
leave the original thesis standing as I first stated it. Scarcity is not, on
its own, a source of innovation. But they do leave a smaller and, I think,
more useful claim standing. It is that innovation has two modes that depend on
each other.

The first mode runs on abundance. It is the mode of research laboratories,
semiconductor fabs, venture capital and national grids, and it produces
\emph{capabilities}: things that did not exist before, such as the solar cell,
the LED, the mobile phone, the intraocular lens and the antibiotic. It is
expensive, slow and concentrated in a few places, and it tends to produce
solutions that assume the rest of the system it grew up in.

The second mode runs on scarcity. It takes those capabilities and remixes them
into forms that work without the surrounding system. It produces fewer
breakthroughs and more \emph{reconfigurations}: a solar cell and an LED in a
box with a payment plan, a phone used as a bank, a surgical procedure turned
into an assembly line. It is cheap, fast and dispersed, and its products tend
to reach more people per dollar, because reaching people without the system
was the whole point.

\begin{keyidea}
Abundance creates capabilities; scarcity remixes them into forms that reach
more people. Neither mode does the other's job, and the second cannot run
without the first.
\end{keyidea}

Seen this way, the hidden abundance in \cref{fig:hidden-abundance} is not an
embarrassment for the argument but part of its structure. The off-grid solar
kit is the scarcity mode working on materials the abundance mode had already
made cheap. M-Pesa is the scarcity mode working on a mobile network that a
multinational had already built. Aravind is the scarcity mode working on a
surgical technique and a lens that rich-country medicine had already
perfected. The reverse flow, when it happens, closes the loop: a reconfigured
product comes back and gives the abundant system a cheaper option it had not
thought to build.

\subsection{When does scarcity produce innovation?}

The cases in this paper are not a sample, so I cannot test conditions
statistically. But reading across them, and across the cases where scarcity
plainly did \emph{not} produce anything, four conditions keep recurring. I
offer them as hypotheses rather than findings.

\textbf{Basic tools are available.} Every success here was built from
components that already existed and were already cheap: bottles and bleach,
mobile phones, LEDs, solar cells, the techniques of cataract surgery. Where
the basic tools themselves are missing, there is nothing to remix. This is the
abundance mode's contribution, and its absence is probably the single most
common reason scarcity produces nothing.

\textbf{Need is dense and shared.} The successful cases served needs that many
people had at once, in the same place: light, sending money, cataracts. Dense,
shared need means that a solution found once can be used many times, and that
people can learn it from their neighbours. Scattered or idiosyncratic needs
do not offer that.

\textbf{There is room to experiment.} This has two parts. The people involved
need enough slack, in money, time and security, to try something that might
fail, which is the lesson of the inverted U in \cref{sec:psychology}. And the
rules need to leave space for unfamiliar models: a regulator willing to let
a phone company hold customers' money, a health system that does not prohibit
task-shifting. Severe deprivation removes the first kind of room; heavy or
unpredictable regulation removes the second.

\textbf{There is a path to scale.} Almost every case that reached millions of
people had some outside source of capital, distribution or credibility:
Vodafone and a donor grant for M-Pesa; domestic manufacturing of supplies and a
training programme that has carried its methods to other hospitals for
Aravind; global supply chains for off-grid solar. Local invention without a path to scale stays local, and
the bottle light, which needed almost no capital, spread largely because a
foundation decided to spread it.

Put together, these conditions describe a particular kind of scarcity. It is
the absence of a \emph{system} in a place where people still have tools,
slack and connections to a wider world. That is a narrower thing than poverty,
and the distinction matters. A policy that wants more of the scarcity mode's
benefits should not aim for more scarcity. It should aim for the conditions
under which scarcity, where it already exists, can become something more than
going without: cheap access to basic components, room to experiment, and
patient capital for the ideas that work.
